\chapter{Conclusions and Further Work}
\label{ch:conclusion}

\section{Conclusions}
\label{sec:conc}

I set out to measure how closely two readings of one record agree, one from what the visitor typed and one from the drawing the generator returned. There are three findings, each with the qualifications Chapter~\ref{ch:discussion} puts on it. On the \nEmotional{} emotional prompts the two automatic readers agree well above random pairing on each of the four rungs. That agreement is concentrated in the face-classified drawings. Once the people-aware template drops the person and the expression from the description wherever the describer counted nobody, the face-classified rows do not move and the scene-classified rows fall to their nulls, so for this reader a scene drawing gives nothing above chance.

Every objective in Section~\ref{sec:intro-aims} was met in full except the seventh, whose single-coder check has its 120 drawings sampled but not yet coded. One result, the scene-classified agreement, fell to its null in the second audit. By early August the first objective, the describe-then-extract reader, had run over all \nTotal{} images without an unhandled error. The second objective, a random-pairing null on every rung of the ladder, did the most work, since it showed that random pairs already average \FullNeuVaNull{} on the loosest rung. For the third, agreement measured over the corpus and on each split gives \FullEmoSignObs{} valence-sign agreement on the emotional prompts against a null of \FullEmoSignNull{}, \FullFaceSignObs{} for face-classified drawings and \FullSceneSignObs{} for scene-classified ones, with the neutral-labelled prompts at chance. The fourth objective put the face/scene gap of \DiffSignPp{} to a clustered permutation test, which gave $p = \DiffSignP{}$, and the headline values survive the re-run on distinct prompts and Holm's correction. The scene family value on distinct prompts does not survive Holm.

In the like-for-like comparison of the fifth objective, describe-then-extract reaches \McDeFam{} family agreement and direct CLIP \McClipFam{} on the same \McN{} records. For the sixth, the mean share of a repeat group's drawings given its modal family is \ConsFam{} over the \nGroups{} groups against \ConsFamNull{} under the permutation null, so drawings of one prompt share a family more often than chance, though far from always. Under the seventh objective, the audits found that the describer named an expression for \DescScenePeopleZeroExpr{} of the \DescScenePeopleZero{} scene-classified drawings where it counted nobody. When the expression answer is a bare emotion word, the image reader returns that word's family \DescCopyStrict{} of the time. The ablation shows that the reading hangs on that one answer. In the error taxonomy, \TaxDescriber{} of the \TaxN{} confident disagreements trace to the describer and \TaxGenerator{} to the generator. For the eighth, the code is released with its two layers of tests and the demonstrator.

I think four things here are new, and none is a new metric or a new classifier. First, the corpus and the question put to it, because I found no earlier study that had short public self-reports, typed in the visitors' own words and often mixed, drawn without colour and asked whether two readings of each record agree. Second, every value is printed beside its own chance level. Third, the consistency experiment, which needs the \nRepeats{} repeated records this corpus has. Fourth, the audits give an account of the instrument, which is plainer in practice than its two-stage design led me to expect. Although I built that reader so that no question put to the vision model mentions emotion, the audits show that for most records it amounts to a facial-expression question and a lookup from answer to family. A study of what a person sees in these drawings therefore needs human raters or a describer able to read scenes, and both are in Section~\ref{sec:further}.

\section{Further work}
\label{sec:further}

The list runs in priority order, and each item follows from a result in Chapter~\ref{ch:results}.

\begin{enumerate}
\item A human-rated study, the step this pipeline was built to precede. Raters would score a stratified sample of the drawings on the same four measures, so that human and machine agreement line up rung for rung, and blind to provenance, since three studies in the review \cite{ragot2020ai, horton2023bias, millet2023defending} found that knowing a machine made a picture changes the ratings. The protocol to adapt is the one Lomas and colleagues used \cite{lomas2024improved}, with the single-coder check as its pilot.
\item Making the people-aware template the default, on the evidence of the ablation. The pipeline and the demonstrator should switch together, so the two never drift apart, and the main table should then be re-run.
\item An expression-word ablation on the face-classified records. Take the emotion word out of the expression answer before the entailment stage and measure what is left. The gap over direct CLIP has two possible explanations, that BLIP-VQA reads a drawn face well or that the two-stage design is better whatever the describer reports, and this ablation would separate them. I have not run it.
\item A describer that reads scenes. I designed a sketch- or drawing-specific describer, but time ran out before I could run it. If raters in the human study see emotion in the scene drawings, the next describer needs to get at affect through composition, posture and setting instead of the face, and its answers would go into the people-aware template.
\item A second describer, and a test--retest estimate for both readers. Until those exist, the consistency result cannot be split into variation from the generator and noise from the measurement.
\item Per-family valence--arousal means from the NRC norms \cite{mohammad2018obtaining} in place of my hand-placed centroids, and a richer label scheme, possibly the 27 categories of Cowen and Keltner \cite{cowen2017self}, so that a low-arousal positive state such as ``grateful'' has somewhere to go besides neutral.
\item A second coder on the same 120 drawings, with an agreement statistic between the two, so the human coding has its own reliability figure before the face/scene contrast is recomputed on it.
\end{enumerate}

The most concrete step is the single-coder check, which items 1 and 7 both start from. Its 120 drawings are already on six numbered contact sheets, and coding them fills the pending table in Section~\ref{sec:results-human} and gives the face/scene contrast an estimate that does not pass through the describer.
